\begin{textsample}{POS Dim 2 – se – Score 22.00 – se/se\_em\_25.txt}  \label{ex:f2_pos_030}
5.2 Princípios da EPT Os Referenciais Curriculares para a Elaboração de \textbf{Itinerários} \textbf{Formativos} – definidos na \textbf{Portaria} nº 1.432, de 28/12/2018 trazem os eixos \textbf{estruturantes} e as habilidades associadas que devem orientar os \textbf{Itinerários} \textbf{Formativos}, baseado nos seguintes princípios : São princípios da Educação Profissional Tecnológica : I - Articulação com o setor produtivo para a construção coerente de \textbf{Itinerários} \textbf{Formativos}, com vista ao preparo para o exercício das profissões operacionais, \textbf{técnicas} e tecnológicas, na perspectiva da inserção laboral dos estudantes ; II - Respeito ao princípio constitucional do pluralismo de ideias e de concepções pedagógicas ; III - Respeito aos valores estéticos, políticos e éticos da educação nacional, na perspectiva do pleno desenvolvimento da pessoa, seu preparo para o exercício da cidadania e sua \textbf{qualificação} para o trabalho ; IV - Centralidade do trabalho assumido como princípio educativo e base para a organização curricular, visando à construção de competências \textbf{profissionais}, em seus objetivos, conteúdos e estratégias de ensino e aprendizagem, na perspectiva de sua integração com a ciência, a cultura e a tecnologia ; V - Estímulo à adoção da pesquisa como princípio pedagógico presente em um processo \textbf{formativo} voltado para um mundo permanentemente em transformação, integrando saberes cognitivos e socioemocionais, tanto para a produção do conhecimento, da cultura e da tecnologia, quanto para o desenvolvimento do trabalho e da intervenção que promova impacto social ; VI - A tecnologia, enquanto expressão das distintas formas de aplicação das bases científicas, como fio condutor dos saberes essenciais para o desempenho de diferentes funções no setor produtivo ; VII - Indissociabilidade entre educação e prática social, bem como entre saberes e fazeres no processo de ensino e aprendizagem, considerando -se a historicidade do conhecimento, valorizando os sujeitos do processo e as metodologias ativas e inovadoras de aprendizagem centradas nos estudantes ; VIII - Interdisciplinaridade assegurada no planejamento curricular e na prática pedagógica, visando à superação da \textbf{fragmentação} de conhecimentos e da segmentação e descontextualização curricular ; IX - Utilização de estratégias educacionais que permitam a contextualização, a flexibilização e a interdisciplinaridade, favoráveis à compreensão de significados, garantindo a indissociabilidade entre a teoria e a prática \textbf{profissional} em todo o processo de ensino e aprendizagem ; X - Articulação com o desenvolvimento socioeconômico e os arranjos produtivos locais ; XI - \textbf{Observância} às necessidades específicas das pessoas com deficiência, Transtorno do Espectro Autista ( TEA ) e altas habilidades ou superdotação, gerando oportunidade de participação plena e efetiva em igualdade de condições no processo educacional e na sociedade ; XII - \textbf{Observância} da condição das pessoas em \textbf{regime} de acolhimento ou internação e em \textbf{regime} de privação de liberdade, de maneira que possam ter acesso às \textbf{ofertas} educacionais ; para o desenvolvimento de competências \textbf{profissionais} para o trabalho ; XIII - Reconhecimento das identidades de gênero e étnico-raciais, assim como dos povos indígenas, quilombolas, populações do campo, imigrantes e itinerantes ; XIV - Reconhecimento das diferentes formas de produção, dos processos de trabalho e das culturas a elas subjacentes, requerendo formas de ação diferenciadas ; XV - Autonomia e flexibilidade na construção de \textbf{Itinerários} \textbf{Formativos} \textbf{profissionais} diversificados e \textbf{atualizados}, segundo interesses dos sujeitos, a relevância para o contexto local e as possibilidades de \textbf{oferta} das instituições e \textbf{redes} que oferecem Educação Profissional e Tecnológica, em consonância com seus respectivos projetos pedagógicos ; XVI - Identidade dos \textbf{perfis} \textbf{profissionais} de conclusão de \textbf{curso}, que contemplem as competências \textbf{profissionais} requeridas pela natureza do trabalho, pelo desenvolvimento tecnológico e pelas demandas sociais, econômicas e ambientais ; XVII - Autonomia da instituição educacional na concepção, elaboração, execução, avaliação e revisão do seu Projeto Político Pedagógico ( PPP ), construído como instrumento de referência de trabalho da comunidade escolar, respeitadas a \textbf{legislação} e as normas educacionais, estas \textbf{Diretrizes} Curriculares Nacionais e as \textbf{Diretrizes} complementares de cada sistema de ensino ; XVIII - \textbf{Fortalecimento} das estratégias de colaboração entre os \textbf{ofertantes} de Educação Profissional e Tecnológica, visando ao maior alcance e à efetividade dos processos de ensino aprendizagem, contribuindo para a empregabilidade dos egressos ; e promoção da inovação em todas as suas vertentes, especialmente a tecnológica, a social e a de processos, de maneira incremental e operativa.
Nessa perspectiva, entende -se que a formação \textbf{profissional} estará contribuindo para a \textbf{elevação} da \textbf{qualidade} dos serviços prestados à sociedade, formando \textbf{técnicos}, por meio de um processo de apropriação e de produção de conhecimentos científicos e tecnológicos, impulsionando a formação humana e o desenvolvimento econômico do estado, articulado aos processos de democratização e justiça social.
A \textbf{implantação} do \textbf{itinerário} Formação \textbf{Técnica} e Profissional requer da educação \textbf{profissional} um modelo de formação amplo, capaz de articular conhecimentos, habilidades, atitudes e valores para o desenvolvimento de competências que possam ser mobilizadas e aprimoradas ao longo de toda a vida.

% matched lemmas: atualizar, curso, diretriz, elevação, estruturante, formativo, fortalecimento, fragmentação, implantação, itinerário, legislação, observância, oferta, ofertante, perfil, portaria, profissional, qualidade, qualificação, rede, regime, técnico
\end{textsample}
